\section{Introduction: the function remains, its human executor does not}
\label{sec:intro}

An enterprise does not need ten buyers, ten contract lawyers, thirty architects, twenty
security architects, and twenty-five production engineers because those numbers are part of its
governance. It needs a particular volume of supplier terms negotiated, clauses checked, architectures
examined, risks assessed, recoveries tested, budgets validated, and decisions recorded. Those people
are the present implementation of that function. An enterprise may still require a contract review, an
architecture review, a security review, or a readiness check after it no longer requires a lawyer, an
architect, or an engineer to perform each one. The institutional function and its human implementation are
different objects.

The relevant replacement is not a model writing a convincing report. It is a system receiving a
dossier, obtaining evidence, investigating exceptions, reaching a defensible decision, and making that
decision effective under a valid mandate. Drafting assistance is not that replacement. Executing the whole work is, and the population
historically needed to run that gate then contracts.

\textbf{The thesis of this paper, the law of gate substitution, is that every gate of a Digital Governance
Framework (DGF) becomes software, and that the workforce executing those gates is reduced
accordingly.} Exceptions, verification, arbitration, and final commitment are included. A job title is not a
permanent barrier to substitution. Neither is describing a decision as judgment: the question is
whether its information, actions, and accepted outcomes can be operationalized. Whatever the role, the
input is a set of documents an agent can read and the output is an artifact an agent can write
(Section~\ref{sec:roles}). We add, as a separate comparative hypothesis, that the DGF will be the first expert control function to
be completely executed by agents (Section~\ref{sec:candidate}); the law does not depend on that
ranking.

The thesis is also what the author observes in practice while operating such frameworks: some automation programs for governance are already planned and justified in terms of fewer posts and
fewer full-time equivalents (FTE), not only faster reviews. The paper takes that objective seriously
and asks what it requires.

The path to that endpoint is not a straight line. A company can automate most of its routine reviews
and end up needing \emph{more} human work: exceptions take longer, outputs require verification, and
maintaining the agents adds an operating burden. Another route on the same platform can save
substantial effort. In the synthetic example of this paper, nine specialist populations supply 140
workload-equivalent FTE to a DGF. At the same case volumes, the same exception fractions, and the
same selection of cases, alternative operating assumptions leave 43.52 or 168.92 FTE of human
workload. Coverage does not identify even the \emph{direction} of the workload change during the
transition. The instrument that explains the gap fits on one line:
\[
\rho=\frac{\phi\eta+(1-\phi)\mu+r+b_A}{1+b_H},
\]
the ratio of human work after deployment to human work before it. The share of baseline labor held by
the exceptions ($\phi$) is multiplied by what deployment does to their effort ($\eta$); the rest is
multiplied by the review still required ($\mu$); rework ($r$) and upkeep ($b_A$) are added
(Section~\ref{sec:frontier}). A reader who rejects the prediction of this paper can still use this
account to evaluate a project. Section~\ref{sec:routes} shows the same fact for two routes of one platform.

The paper therefore does three things. It \emph{argues} for the destination from five premises:
contracts already delimit the work; measured agent capabilities are extending to longer tasks; bounded
gates already execute as software; authority can be delegated through standing rules; and common
platforms can amortize the cost of reliable execution. It supplies the \emph{instrument} that measures
the path: a complete human-labor account, a selection identity explaining why a small exception queue
can hold most of the work, a break-even frontier, and fragility thresholds. And it gives the workforce prediction a
\emph{deadline}: 80\% fewer people working in the DGF ecosystem by 20 September 2033 than in 2026,
with failure criteria that cannot be escaped by renaming remaining workers supervisors. Complete
automation is the longer-term conjecture, without a fixed date. Workforce displacement follows from successful
substitution; it is the consequence of the thesis, not the evidence for it.

\medskip
\noindent\fbox{\begin{minipage}{\dimexpr\linewidth-2\fboxsep-2\fboxrule\relax}\small
\textbf{Scope and status.} This paper states a thesis and the reasons for it; it is not an empirical
demonstration. No autonomous DGF deployment or enterprise staffing survey is reported. All staffing
numbers and calendar scenarios are synthetic calculations. Product documentation establishes described
functionality, not its prevalence. Legal provisions constrain the argument and require
deployment-specific review. Workload-equivalent FTE measures required effort for a fixed governed
output; actual payroll also depends on staffing adjustment, demand, and new tasks. Gates are interchangeable: each enterprise
composes its own DGF, and the gates, routes, and orders used here are an arbitrary selection taken as
examples. The universal endpoint is a conjecture; the separate 80\% workforce-reduction hypothesis for 2033 is deliberately refutable.
\end{minipage}}
